
\subsection{含幺情形：无滤过版本}
\label{subsec:4-5-3}

本小节证明定理 4.2.45a 和 4.2.45b 的无滤过版本。对定理
4.2.45a 的无滤过版本，证明与定理 4.2.45 (1) 的证明平行；
主要差别在于，须以以下定理 4.5.1a 代替定理 4.5.1。
设 $\overline{\mathfrak f}\colon\overline C_1\to\overline C_2$
为含幺 $A_K$ 同态。

回忆
\[
 B_k\overline C_{1,\mathrm{red}}[1]
 =
 \underbrace{\overline C_1[1]/R[e]\otimes\cdots\otimes
 \overline C_1[1]/R[e]}_{k\text{ 个因子}}.
\]
由于 $\overline{\mathfrak m}_1e=0$，
$\widehat{\overline{\mathfrak m}}_1$ 在
$B_k\overline C_{1,\mathrm{red}}[1]$ 上诱导算子，
仍记作 $\widehat{\overline{\mathfrak m}}_1$。
因此它与
$\overline{\mathfrak m}_1\colon\overline C_2[1]\to
\overline C_2[1]$
在
$\operatorname{Hom}(B_k\overline C_{1,\mathrm{red}}[1],
\overline C_2[1])$
上诱导微分 $\delta_1$。

\begin{theorem}{4.5.1a}\label{thm:4-5-1a}
存在次数为 $1$ 的
\[
 o^{\mathrm{red}}_{K+1}(\overline{\mathfrak f})
 \in\operatorname{Hom}
 (B_{K+1}\overline C_{1,\mathrm{red}}[1],
  \overline C_2[1]),
\]
使得下列结论成立：
\begin{enumerate}[label=\textup{(4.5.2a.\arabic*)},leftmargin=6.1em]
 \item
 $\delta_1(o^{\mathrm{red}}_{K+1}
 (\overline{\mathfrak f}))=0$。
 \item
 $[o^{\mathrm{red}}_{K+1}(\overline{\mathfrak f})]=0$
 作为 $\delta_1$-上同调类，当且仅当存在延拓
 $\overline{\mathfrak f}$ 的 $A_{K+1}$ 同态。
 \item
 若 $\overline{\mathfrak f}'$ 与
 $\overline{\mathfrak f}$ 是 $A_K$ 同伦的，则
 $[o^{\mathrm{red}}_{K+1}(\overline{\mathfrak f})]
 =[o^{\mathrm{red}}_{K+1}(\overline{\mathfrak f}')]$
 作为 $\delta_1$-上同调类。
 \item
 若 $\overline{\mathfrak g}\colon
 \overline C'_1\to\overline C_1$ 和
 $\overline{\mathfrak g}'\colon
 \overline C_2\to\overline C'_2$
 都是含幺 $A_{K+1}$ 同态，则
 \[
 \begin{aligned}
 &[o^{\mathrm{red}}_{K+1}
   (\overline{\mathfrak g}'\circ
    \overline{\mathfrak f}\circ\overline{\mathfrak g})]\\
 &\qquad=
 (\overline{\mathfrak g}'_1)_*
 \circ[o^{\mathrm{red}}_{K+1}(\overline{\mathfrak f})]
 \circ(\overline{\mathfrak g}_1\otimes\cdots
 \otimes\overline{\mathfrak g}_1)_*
 \end{aligned}
 \]
 在
 $H(\operatorname{Hom}
 (B_{K+1}\overline C'_{1,\mathrm{red}}[1],
 \overline C'_2[1]),\delta_1)$
 中成立。\translatornote{原书此处的上同调空间把带撇的源、靶
 写成无撇对象；但左端复合同态的源、靶分别是
 $\overline C'_1$、$\overline C'_2$，译文依此定型。}
 这里，
 \[
 \overline{\mathfrak g}_1\otimes\cdots
 \otimes\overline{\mathfrak g}_1\colon
 B_{K+1}\overline C_{1}^{\prime\,\mathrm{red}}[1]
 \longrightarrow B_{K+1}\overline C_{1}^{\mathrm{red}}[1]
 \]
 由 $\overline{\mathfrak g}_1\colon
 \overline C_{1}^{\prime\,\mathrm{red}}[1]\to
 \overline C_{1}^{\mathrm{red}}[1]$ 诱导。
\end{enumerate}
\end{theorem}

\begin{proof}
用引理 4.4.30（或其 $A_K$ 对应版本），对
$k=1,\ldots,K$ 配以 $R$-模同态
\[
 \widehat{\overline{\mathfrak f}}_{k,\mathrm{red}}\colon
 B_k^{\mathrm{red}}\overline C_1[1]
 \longrightarrow\overline C_2[1].
\]
令 $\widehat{\overline{\mathfrak f}}_{K+1,\mathrm{red}}=0$，
它们诱导
\[
 \widehat{\overline{\mathfrak f}}^{\mathrm{red}}_{1\cdots K+1}
 \in\operatorname{Hom}
 (B_{1\cdots K+1,\mathrm{red}}\overline C_1[1],
  B_{1\cdots K+1}\overline C_2[1]).
\]
于是考虑
\[
 \begin{aligned}
 o^{\mathrm{red},0}_{K+1}(\overline{\mathfrak f})
 &=
 \widehat d_{1\cdots K+1}\circ
 \widehat{\overline{\mathfrak f}}_{1\cdots K+1}
 -
 \widehat{\overline{\mathfrak f}}_{1\cdots K+1}
 \circ\widehat d_{1\cdots K+1}\\
 &\in\operatorname{Hom}
 (B_{K+1}^{\mathrm{red}}\overline C_1[1],
  \overline C_2[1]).
 \end{aligned}
\]
仿引理 4.4.35 的证明可得
$o^{\mathrm{red},0}_{K+1}\circ s_{\mathrm{left}}
=o^{\mathrm{red},0}_{K+1}\circ s_{\mathrm{right}}=0$。
因此得到
$o^{\mathrm{red}}_{K+1}(\overline{\mathfrak f})
\in\operatorname{Hom}
(B_{K+1}\overline C_{1,\mathrm{red}}[1],
\overline C_2[1])$。照定理 4.5.1 的证明逐项验证，
即有以上各性质。
\end{proof}

用定理 4.5.1a 代替定理 4.5.1，其余证明照旧，
便得到定理 4.2.45a 的无滤过版本。

再考察定理 4.2.45b 的无滤过版本。从一个同伦含幺
$A_\infty$ 同态 $\overline C_1\to\overline C_2$ 出发；
按定义，它是一个含幺 $A_\infty$ 同态
$\overline{\mathfrak f}^{+}\colon
\overline C_1^{+}\to\overline C_2^{+}$，并满足
\begin{equation}
 \overline{\mathfrak f}^{+}_k
 \bigl(B_k(\overline C_1[1])\bigr)
 \subseteq\overline C_2[1].
\tag{4.5.11}\label{eq:4-5-11}
\end{equation}
需要找到同伦逆
$\overline{\mathfrak g}^{+}\colon
\overline C_2^{+}[1]\to\overline C_1^{+}[1]$，使
\begin{equation}
 \overline{\mathfrak g}^{+}_k
 \bigl(B_k(\overline C_2[1])\bigr)
 \subseteq\overline C_1[1].
\tag{4.5.12}\label{eq:4-5-12}
\end{equation}
故要对
$\overline{\mathfrak f}^{+}$
应用定理 4.2.45a，还需补充论证。

取 $[0,1]\times\overline C_1^{+}$ 的含幺模型
$\overline{\mathfrak C}_1^{(+)}$，置
\begin{equation}
 \overline{\mathfrak C}_1
 =
 (\overline{\mathrm{Eval}}_{s=0})_1^{-1}
 (\overline C_1)
 \cap
 (\overline{\mathrm{Eval}}_{s=1})_1^{-1}
 (\overline C_1).
\tag{4.5.13}\label{eq:4-5-13}
\end{equation}
由于 $k\ne1$ 时
$(\overline{\mathrm{Eval}}_{s=s_0})_k=0$，
$\overline{\mathfrak C}_1$ 是
$A_\infty$ 子代数；还要注意，
$\overline C_1$ 是 $\overline C_1^+$ 的
$A_\infty$ 子代数。容易看出，
$\overline{\mathfrak C}_1$ 是
$[0,1]\times\overline C_1$ 的模型。

现在对 $K$ 归纳证明以下命题。
\begin{proposition}{4.5.6b}\label{prop:4-5-6b}
设含幺弱同伦等价
$\overline{\mathfrak f}^{+}\colon
\overline C_1^+\to\overline C_2^+$
满足 \eqref{eq:4-5-11}，
含幺 $A_K$ 同态
$\overline{\mathfrak g}^{+(K)}\colon
\overline C_2^+\to\overline C_1^+$
满足 \eqref{eq:4-5-12}。又设
$\overline{\mathfrak h}^{+(K)}\colon
\overline C_1^+\to\overline{\mathfrak C}_1^{(+)}$
为从恒等同态到
$\overline{\mathfrak g}^{+(K)}
\circ\overline{\mathfrak f}^{+}$
的含幺 $A_K$ 同伦，并满足
\begin{equation}
 \overline{\mathfrak h}^{+}_k
 \bigl(B_k(\overline C_1[1])\bigr)
 \subseteq\overline{\mathfrak C}_1[1].
\tag{4.5.14}\label{eq:4-5-14}
\end{equation}
则 $\overline{\mathfrak g}^{+(K)}$
可延拓为满足 \eqref{eq:4-5-12} 的
含幺 $A_{K+1}$ 同态
$\overline{\mathfrak g}^{+(K+1)}$，
而 $\overline{\mathfrak h}^{+(K)}$
可延拓为满足 \eqref{eq:4-5-14} 的
含幺 $A_{K+1}$ 同伦
$\overline{\mathfrak h}^{+(K+1)}$，
从恒等同态到
$\overline{\mathfrak g}^{+(K+1)}
\circ\overline{\mathfrak f}^{+}$。
\end{proposition}
\begin{proof}
注意以下两点后，证明就与定理 4.5.1a 的证明平行。
首先，$A_\infty$ 代数
$\overline C_i^+$ 有典范分裂
$\overline C_i^+
=\overline C_i^{+\mathrm{red}}\oplus Re$，
其中
$\overline C_i^{+\mathrm{red}}
=\overline C_i\oplus Rf$。
其次，为保持 \eqref{eq:4-5-12}
和 \eqref{eq:4-5-14}，须稍改定理
4.5.1a 的障碍理论。例如，将
$\overline{\mathfrak g}^{+(K)}$
延拓为仍满足 \eqref{eq:4-5-12} 的
含幺 $A_{K+1}$ 同态，其障碍须位于受限空间
\translatornote{原书这两组障碍空间沿用
$\overline C_1\to\overline C_2$ 的方向，
与本命题的
$\overline{\mathfrak g}^{+(K)}
\colon\overline C_2^+\to\overline C_1^+$
不符；此处依同态方向及条件 (4.5.12) 定型。}
\[
 \operatorname{Hom}
 (B_{K+1}\overline C_2^{+\mathrm{red}}[1],
 \overline C_1[1]),
 \quad\text{而非一般的}\quad
 \operatorname{Hom}
 (B_{K+1}\overline C_2^{+\mathrm{red}}[1],
 \overline C_1^{+\mathrm{red}}[1]).
\]
其余证明同样进行。由此命题及定理 4.2.45b 的无滤过版本得证。
\end{proof}

\subsection{将模 $T^{\lambda_i}$ 的滤过 $A_\infty$ 同态延拓至模 $T^{\lambda_{i+1}}$}
\label{subsec:4-5-4}

下面着手证明定理 4.2.45 (2)，为此对能量滤过归纳。
令 $G\subset\mathbb R_{\ge0}\times2\mathbb Z$
满足条件 3.1.6，并假定所论的所有滤过 $A_\infty$
代数及滤过 $A_\infty$ 同态均为 $G$-有隙。
与第 4.4.6 小节一样，可将 $G$ 到
$\mathbb R_{\ge0}$ 的投影像列为
$\operatorname{Spec}(G)=\{\lambda_0,\lambda_1,\ldots\}$，
使 $\lambda_i<\lambda_{i+1}$ 且
$\lim_{i\to\infty}\lambda_i=\infty$。
令
\[
 \mathfrak f_k\colon B_k C_1[1]\to C_2[1],
 \qquad k=0,1,2,\ldots,
\]
是一列满足 (3.2.27) 的次数 $0$ 同态。
它们如 (3.2.28) 诱导余代数同态
$\widehat{\mathfrak f}\colon
\widehat B C_1[1]\to\widehat B C_2[1]$。

\begin{definition}{4.5.15}\label{def:4-5-15}
若
\[
 \widehat d\circ\widehat{\mathfrak f}
 -\widehat{\mathfrak f}\circ\widehat d
 \equiv0\pmod{T^{\lambda_i}},
\]
则称 $\mathfrak f$ 为\emph{模 $T^{\lambda_i}$
的滤过 $A_\infty$ 同态}\liftIndex{I327}。这里 $i=1,2,\ldots$。
\end{definition}

设 $\mathfrak f,\mathfrak g$ 是从 $C_1$ 到 $C_2$
的模 $T^{\lambda_i}$ 的滤过 $A_\infty$ 同态，
并取 $[0,1]\times C_2$ 的模型 $\mathfrak C_2$。
若存在模 $T^{\lambda_i}$ 的滤过 $A_\infty$ 同态
$\mathfrak h\colon C_1\to\mathfrak C_2$，使
\[
 \mathrm{Eval}_{s=0}\circ\mathfrak h
 \equiv\mathfrak f\pmod{T^{\lambda_i}},
 \qquad
 \mathrm{Eval}_{s=1}\circ\mathfrak h
 \equiv\mathfrak g\pmod{T^{\lambda_i}},
\]
便称二者\emph{模 $T^{\lambda_i}$ 同伦}。
\translatornote{原书此处把 $\mathfrak h$ 的靶写成
$C_2$；评价同态作用的对象必须是
$[0,1]\times C_2$ 的模型 $\mathfrak C_2$。}

模 $T^{\lambda_i}$ 的滤过 $A_\infty$ 同态的复合，
可按显然方式定义。还须注意，从 $C_1$ 到 $C_2$
的这种同态 $\mathfrak f$，在 $R$-约化后
诱导 $A_\infty$ 同态
$\overline{\mathfrak f}\colon
\overline C_1\to\overline C_2$。

现在可以开始证明定理 4.2.45 (2)。
\begin{statement}{定义兼命题}{4.5.16}\label{defprop:4-5-16}
设 $\mathfrak f^{(i)}$ 是定义 4.5.15
意义下的模 $T^{\lambda_i}$ 的滤过
$A_\infty$ 同态。则存在 Hochschild 上同调类\liftIndex{I352}
\[
 [o_{\lambda_i}(\mathfrak f^{(i)})]
 \in HH(\overline C_1,\overline C_2;
 \overline{\mathfrak f}^{(i)}),
\]
满足：
\begin{enumerate}[label=\textup{(4.5.17.\arabic*)},leftmargin=5.7em]
 \item
 $[o_{\lambda_i}(\mathfrak f^{(i)})]=0$
 当且仅当存在 $\mathfrak f^{(i)\prime}$，使
 \begin{enumerate}[label=\textup{(4.5.17.1.\arabic*)},leftmargin=7.4em]
  \item
  $\mathfrak f^{(i)\prime}$ 是模
  $T^{\lambda_{i+1}}$ 的滤过 $A_\infty$ 同态；
  \item
  $\mathfrak f^{(i)\prime}\equiv
  \mathfrak f^{(i)}\pmod{T^{\lambda_i}}$。
 \end{enumerate}
 \item
 若 $\mathfrak f^{(i)\prime}$ 也是模
 $T^{\lambda_i}$ 的滤过 $A_\infty$ 同态，
 且与 $\mathfrak f^{(i)}$ 模 $T^{\lambda_i}$
 同伦，则
 \[
 [o_{\lambda_i}(\mathfrak f^{(i)})]
 =[o_{\lambda_i}(\mathfrak f^{(i)\prime})].
 \]
 \item
 若 $\mathfrak g\colon C'_1\to C_1$、
 $\mathfrak g'\colon C_2\to C'_2$
 是模 $T^{\lambda_{i+1}}$ 的滤过
 $A_\infty$ 同态，则
 \[
 [o_{\lambda_i}(\mathfrak g'\circ
 \mathfrak f^{(i)}\circ\mathfrak g)]
 =(\overline{\mathfrak g},
 \overline{\mathfrak g}')_*
 [o_{\lambda_i}(\mathfrak f^{(i)})].
 \]
\end{enumerate}
\end{statement}
\liftIndex{I170}

\begin{remark}{4.5.18}\label{rem:4-5-18}
须解释 (4.5.17.2) 的含义。由于假设
$\overline{\mathfrak f}^{(i)}$、
$\overline{\mathfrak f}^{(i)\prime}
\colon\overline C_1\to\overline C_2$
由同伦
$\overline{\mathfrak h}\colon
\overline C_1\to\overline{\mathfrak C}_2$
相连，存在典范同构
\[
 HH(\overline C_1,\overline C_2;
 \overline{\mathfrak f}^{(i)})
 \cong
 HH(\overline C_1,\overline C_2;
 \overline{\mathfrak f}^{(i)\prime}).
\]
因此 (4.5.17.2) 的等式是在该同构下理解的。
这一同构源于以下事实：
$(\overline{\mathrm{Eval}}_{s=0})_1$
和 $(\overline{\mathrm{Incl}})_1$ 诱导同构
\[
 HH(\overline C_1,\overline{\mathfrak C}_2;
 \overline{\mathfrak h})
 \cong
 HH(\overline C_1,\overline C_2;
 \overline{\mathfrak f}^{(i)}),
\]
而 $(\overline{\mathrm{Eval}}_{s=1})_1$
和 $(\overline{\mathrm{Incl}})_1$ 诱导同构
\[
 HH(\overline C_1,\overline{\mathfrak C}_2;
 \overline{\mathfrak h})
 \cong
 HH(\overline C_1,\overline C_2;
 \overline{\mathfrak f}^{(i)\prime}).
\]
这由命题 4.4.48 得出。
\end{remark}

\begin{proof}[定义兼命题 4.5.16 的证明]
由假设，存在满足以下条件的障碍项：
\[
 \begin{gathered}
 o_{\lambda_i}(\mathfrak f^{(i)})
 \in\operatorname{Hom}
 (B\overline C_1[1],B\overline C_2[1]),\\
 \widehat d\circ\widehat{\mathfrak f}^{(i)}
 -\widehat{\mathfrak f}^{(i)}\circ\widehat d
 \equiv
 T^{\lambda_i}o_{\lambda_i}(\mathfrak f^{(i)})
 \pmod{T^{\lambda_{i+1}}}.
 \end{gathered}
\]
由于 $\widehat d$ 是余导子，
$\widehat{\mathfrak f}^{(i)}$ 是余代数同态，
$o_{\lambda_i}(\mathfrak f^{(i)})$
也是余导子。而且
\translatornote{原书以下恒等式的第二括号误将交换子写成
$\widehat d\circ\widehat{\mathfrak f}^{(i)}
-\widehat d\circ\widehat{\mathfrak f}^{(i)}$；
此处按闭性所需的交换子恢复。}
\[
 \begin{aligned}
 &\widehat d\bigl(
   \widehat d\circ\widehat{\mathfrak f}^{(i)}
   -\widehat{\mathfrak f}^{(i)}\circ\widehat d\bigr)\\
 &\quad+
 \bigl(\widehat d\circ\widehat{\mathfrak f}^{(i)}
 -\widehat{\mathfrak f}^{(i)}\circ\widehat d\bigr)
 \widehat d=0,
 \end{aligned}
\]
故 $\delta(o_{\lambda_i}(\mathfrak f^{(i)}))=0$。
于是它给出
$[o_{\lambda_i}(\mathfrak f^{(i)})]
\in HH(\overline C_1,\overline C_2;
\overline{\mathfrak f}^{(i)})$。
仿定理 4.5.1 的证明，即可验证
(4.5.17.1)、(4.5.17.2)、(4.5.17.3)。
\end{proof}

\subsection{定理 4.2.45 的证明 II：能量滤过}
\label{subsec:4-5-5}
\liftIndex{I171}

回到定理 4.2.45 (2) 的证明。
沿用上一小节的 $G$ 与 $\lambda_i$。
令 $C_i$ 为 $G$-有隙滤过 $A_\infty$ 代数，
令 $\mathfrak f\colon C_1\to C_2$
为 $G$-有隙滤过 $A_\infty$ 同态。
假定 $\mathfrak f$ 是弱同伦等价，
即由它诱导的无滤过 $A_\infty$ 同态
$\overline{\mathfrak f}\colon
\overline C_1\to\overline C_2$
在 $\overline{\mathfrak m}_1$-上同调上诱导同构。
由已经证明的定理 4.2.45 (1)，
$\overline{\mathfrak f}$ 是同伦等价。

对 $n$ 归纳证明以下命题。
\begin{proposition}{4.5.19}\label{prop:4-5-19}
设 $\mathfrak g^{(n)}\colon C_2\to C_1$
是模 $T^{\lambda_n}$ 的滤过 $A_\infty$ 同态，
而 $\mathfrak h^{(n)}\colon C_1\to\mathfrak C_1$
是模 $T^{\lambda_n}$ 的滤过 $A_\infty$ 同伦，
连接 $\mathrm{id}$ 与
$\mathfrak g^{(n)}\circ\mathfrak f$。
则存在模 $T^{\lambda_{n+1}}$ 的滤过
$A_\infty$ 同态
$\mathfrak g^{(n+1)}\colon C_2\to C_1$，
以及模 $T^{\lambda_{n+1}}$ 的滤过
$A_\infty$ 同伦
$\mathfrak h^{(n+1)}\colon C_1\to\mathfrak C_1$，
连接 $\mathrm{id}$ 与
$\mathfrak g^{(n+1)}\circ\mathfrak f$。
此外，
\[
 \mathfrak g^{(n+1)}
 \equiv\mathfrak g^{(n)}\pmod{T^{\lambda_n}},
 \qquad
 \mathfrak h^{(n+1)}
 \equiv\mathfrak h^{(n)}\pmod{T^{\lambda_n}}.
\]
\end{proposition}

\begin{center}
\resizebox{0.96\linewidth}{!}{$
\begin{array}{ccccccc}
 C_1
 &\xleftarrow{\ \mathrm{Eval}_{s=0}\ }&
 \mathfrak C_1
 &\xrightarrow{\ \mathrm{Eval}_{s=1}\ }&
 C_1
 &\xleftarrow{\ \mathfrak g^{(n)}\ }&
 C_2\\[0.5em]
 \Vert&&\underset{\scriptstyle\mathfrak h}{\Big\uparrow}
 &&&&\Vert\\[0.5em]
 C_1
 &\mathrel{\hspace{1.7em}=\hspace{1.7em}}&
 C_1
 &\mathrel{\hspace{1.7em}=\hspace{1.7em}}&
 C_1
 &\xrightarrow{\ \mathfrak f\ }&
 C_2
\end{array}
$}
\par\smallskip\textbf{图式 4.5.2}\label{diag:4-5-2}
\end{center}

\begin{proof}
考察 $o_{\lambda_n}(\mathfrak h^{(n)})$。
由于
$(\overline{\mathrm{Eval}}_{s=0},
\mathrm{id})_*
[o_{\lambda_n}(\mathfrak h^{(n)})]
=o_{\lambda_n}(\mathrm{id})=0$，
由命题 4.4.48 和 (4.5.2.4) 可知
$[o_{\lambda_n}(\mathfrak h^{(n)})]=0$。
事实上还能得到更强的结论
\translatornote{原书式 (4.5.20) 的同伦模型及评价靶
误写为下标 $2$；图式 4.5.2 与本命题都规定
$\mathfrak h^{(n)}$ 的靶是 $\mathfrak C_1$。}。
\begin{equation}
 \begin{aligned}
 \mathfrak K:=\ker\bigl(
 (\overline{\mathrm{Eval}}_{s=0},\mathrm{id})_*
 \colon&
 \operatorname{Der}
 (B\overline C_1[1],
 B\overline{\mathfrak C}_1[1];
 \overline{\mathfrak h})\\
 &\longrightarrow
 \operatorname{Der}
 (B\overline C_1[1],
 B\overline C_1[1];\mathrm{id})
 \bigr).
 \end{aligned}
\tag{4.5.20}\label{eq:4-5-20}
\end{equation}
障碍
$o_{\lambda_n}(\mathfrak h^{(n)})$
落在 $\mathfrak K$ 中。由命题 4.4.48
及 (4.2.3.3)，$\mathfrak K$ 是无环复形。
故存在 $\mathfrak h'_{n+1}\in\mathfrak K$，满足
\[
 -\delta(\mathfrak h'_{n+1})
 =o_{\lambda_n}(\mathfrak h^{(n)}),
 \qquad
 (\overline{\mathrm{Eval}}_{s=0},\mathrm{id})_*
 (\mathfrak h'_{n+1})=0.
\]
于是
$\mathfrak h^{\prime(n+1)}
=\mathfrak h^{(n)}
+T^{\lambda_n}\mathfrak h'_{n+1}$
是模 $T^{\lambda_{n+1}}$ 的滤过
$A_\infty$ 同态。
\translatornote{原书把这里的修正项写成
$T^{\lambda_{n+1}}\mathfrak h'_{n+1}$；
待消去的首个障碍是 $T^{\lambda_n}$ 阶，
故修正须从该阶开始。下文求
$\mathfrak g'_{n+1}$ 时，原书漏掉的负号
也依同一障碍定义恢复。}

又因为
\[
 -\delta\bigl(
 (\overline{\mathrm{Eval}}_{s=1},\mathrm{id})_*
 (\mathfrak h'_{n+1})\bigr)
 =
 (\overline{\mathrm{Eval}}_{s=1},\mathrm{id})_*
 o_{\lambda_n}(\mathfrak h^{(n)})
 =
 o_{\lambda_n}(\mathfrak g^{(n)}\circ\mathfrak f),
\]
且 $\overline{\mathfrak f}$ 在 Hochschild 上同调上
诱导同构（命题 4.4.48），可取
$\mathfrak g'_{n+1}
\in\operatorname{Der}
(B\overline C_2[1],B\overline C_1[1];
\overline{\mathfrak g})$，使
$\delta(\mathfrak g'_{n+1})
=-o_{\lambda_n}(\mathfrak g^{(n)})$。
因而
\[
 \delta\bigl(
 (\overline{\mathrm{Eval}}_{s=1},\mathrm{id})_*
 (\mathfrak h'_{n+1})
 -
 (\mathrm{id},\overline{\mathfrak f})_*
 (\mathfrak g'_{n+1})\bigr)=0.
\]

再由 $\overline{\mathfrak f}$ 在 Hochschild
上同调上诱导同构，可取下列两个修正项，
使其满足随后的等式：
\[
 \Delta\mathfrak g_{n+1}
 \in\operatorname{Der}
 (B\overline C_2[1],B\overline C_1[1];
 \overline{\mathfrak g}),
 \qquad
 \Delta_1\mathfrak h_{n+1}
 \in\operatorname{Der}
 (B\overline C_1[1],B\overline C_1[1];
 \overline{\mathfrak g}\circ\overline{\mathfrak f}),
\]
\[
 \begin{aligned}
 &(\mathrm{id},\overline{\mathfrak f})_*
 (\mathfrak g'_{n+1})
 +(\mathrm{id},\overline{\mathfrak f})_*
 (\Delta\mathfrak g_{n+1})\\
 &\qquad
 -(\overline{\mathrm{Eval}}_{s=1},\mathrm{id})_*
 (\mathfrak h'_{n+1})
 =\delta(\Delta_1\mathfrak h_{n+1}).
 \end{aligned}
\]
\clearpage
根据 (4.2.3.3) 及引理 4.4.43，
又可取满足下列端点条件的修正项
\translatornote{原书此段把两个同伦修正的源误写为
$C_2$，并混入 $K+1$ 与多出的撇号；
它们都须以 $C_1$ 为源，本段归纳阶为 $n+1$。}
\[
 \Delta\mathfrak h_{n+1}
 \in\operatorname{Der}
 (B\overline C_1[1],B\overline{\mathfrak C}_1[1];
 \overline{\mathfrak h}),
\]
\[
 (\overline{\mathrm{Eval}}_{s=1})_*
 (\Delta\mathfrak h_{n+1})
 =\Delta_1\mathfrak h_{n+1},
 \qquad
 (\overline{\mathrm{Eval}}_{s=0})_*
 (\Delta\mathfrak h_{n+1})=0.
\]
最后，定义两个修正系数：
\[
 \mathfrak g_{n+1}
 =\mathfrak g'_{n+1}+\Delta\mathfrak g_{n+1},
 \qquad
 \mathfrak h_{n+1}
 =\mathfrak h'_{n+1}
  +\delta(\Delta\mathfrak h_{n+1}).
\]
以这两个修正系数在 $T^{\lambda_n}$ 阶
修正 $\mathfrak g^{(n)}$ 与
$\mathfrak h^{(n)}$，得到
$\mathfrak g^{(n+1)}$ 和
$\mathfrak h^{(n+1)}$。
容易逐项核对，它们满足命题所需性质。
\end{proof}

现在完成定理 4.2.45 (2) 的证明。
第 4.5.2 小节已证的定理 4.2.45 (1)
表明命题 4.5.19 的假设在 $n=1$ 时成立。
设归纳已得
$\mathfrak g^{(n)}$ 与
$\mathfrak h^{(n)}$。
由于
$\mathfrak g^{(n+1)}
\equiv\mathfrak g^{(n)}
\pmod{T^{\lambda_n}}$、
$\mathfrak h^{(n+1)}
\equiv\mathfrak h^{(n)}
\pmod{T^{\lambda_n}}$，
两序列的极限存在。
于是得到 $\mathfrak g$，
使 $\mathfrak g\circ\mathfrak f$
同伦于恒等同态。
证明 $\mathfrak f\circ\mathfrak g\sim\mathrm{id}$
的方法，与定理 4.2.45 (1) 证明的最后一步相同。
至此定理 4.2.45 得证。

对定理 4.2.45a、4.2.45b 的含幺或同伦含幺版本，
只需修改上述论证。所需修改与
第 4.5.3、4.4.7 小节的方法类似，故略去细节。

\begin{proof}[命题 4.3.16 的证明]
须证明 $\mathfrak f\equiv\mathrm{id}
\pmod{\Lambda^+_{0,\mathrm{nov}}}$。
沿用第 4.3 节证明中的记号。
注意，若将
$(\mathfrak C,\mathfrak m^{\widetilde b})$
的系数约化为 $R$，便得到
$[0,1]\times\overline C$ 的一个模型。
因此，构造
$\mathrm{Eval}_{s=0}\colon
(\mathfrak C,\mathfrak m^{\widetilde b})
\to(C,\mathfrak m^{b_0})$
的同伦逆时，可从
$\overline{\mathrm{Incl}}$ 开始归纳。
从而所得同伦等价
$\mathfrak f\colon
(C,\mathfrak m^{b_0})
\to(C,\mathfrak m^{b_1})$
满足所需的
$\overline{\mathfrak f}
=\overline{\mathrm{Eval}}_{s=1}
\circ\overline{\mathrm{Incl}}
=\mathrm{id}$。
\end{proof}
